\backmatter
\chapter{References}
\begin{thebibliography}{99}

\bibitem{Riemann1859}
B. Riemann,
\emph{Ueber die Anzahl der Primzahlen unter einer gegebenen Grösse},
Monatsberichte der Berliner Akademie, 1859, pp. 671--680.

\bibitem{Shannon1948}
C. E. Shannon,
``A Mathematical Theory of Communication,''
\emph{Bell System Technical Journal}, vol. 27, pp. 379--423 and 623--656, 1948.
DOI: \href{https://doi.org/10.1002/j.1538-7305.1948.tb01338.x}{10.1002/j.1538-7305.1948.tb01338.x} and
\href{https://doi.org/10.1002/j.1538-7305.1948.tb00917.x}{10.1002/j.1538-7305.1948.tb00917.x}.

\bibitem{Titchmarsh1986}
E. C. Titchmarsh, revised by D. R. Heath-Brown,
\emph{The Theory of the Riemann Zeta-Function}, 2nd ed., Oxford University Press, 1986.

\bibitem{Edwards1974}
H. M. Edwards,
\emph{Riemann's Zeta Function}, Academic Press, 1974.

\bibitem{Bombieri2006}
E. Bombieri,
``The Riemann Hypothesis,'' in J. Carlson, A. Jaffe, and A. Wiles, eds.,
\emph{The Millennium Prize Problems}, Clay Mathematics Institute and American Mathematical Society, 2006.

\bibitem{DLMF25}
NIST Digital Library of Mathematical Functions,
Chapter 25, ``Zeta and Related Functions,'' especially Sections 25.4 and 25.10,
\url{https://dlmf.nist.gov/25}, accessed 29 July 2026.

\bibitem{DimitrovXu2016}
D. K. Dimitrov and Y. Xu,
``Wronskians of Fourier and Laplace Transforms,''
\emph{arXiv preprint arXiv:1606.05011}, 2016.
\url{https://arxiv.org/abs/1606.05011}.

\bibitem{AmariNagaoka2000}
S.-I. Amari and H. Nagaoka,
\emph{Methods of Information Geometry},
American Mathematical Society and Oxford University Press, 2000.

\bibitem{Berry1984}
M. V. Berry,
``Quantal Phase Factors Accompanying Adiabatic Changes,''
\emph{Proceedings of the Royal Society of London A}, vol. 392, no. 1802, pp. 45--57, 1984.
DOI: \href{https://doi.org/10.1098/rspa.1984.0023}{10.1098/rspa.1984.0023}.

\bibitem{Simon1983}
B. Simon,
``Holonomy, the Quantum Adiabatic Theorem, and Berry's Phase,''
\emph{Physical Review Letters}, vol. 51, no. 24, pp. 2167--2170, 1983.
DOI: \href{https://doi.org/10.1103/PhysRevLett.51.2167}{10.1103/PhysRevLett.51.2167}.

\bibitem{AharonovBohm1959}
Y. Aharonov and D. Bohm,
``Significance of Electromagnetic Potentials in the Quantum Theory,''
\emph{Physical Review}, vol. 115, no. 3, pp. 485--491, 1959.
DOI: \href{https://doi.org/10.1103/PhysRev.115.485}{10.1103/PhysRev.115.485}.

\bibitem{Hubble1929}
E. Hubble,
``A Relation between Distance and Radial Velocity among Extra-Galactic Nebulae,''
\emph{Proceedings of the National Academy of Sciences}, vol. 15, no. 3, pp. 168--173, 1929.
DOI: \href{https://doi.org/10.1073/pnas.15.3.168}{10.1073/pnas.15.3.168}.

\bibitem{BerryKeating1999}
M. V. Berry and J. P. Keating,
``The Riemann Zeros and Eigenvalue Asymptotics,''
\emph{SIAM Review}, vol. 41, no. 2, pp. 236--266, 1999.
DOI: \href{https://doi.org/10.1137/S0036144598347497}{10.1137/S0036144598347497}.

\bibitem{Connes1999}
A. Connes,
``Trace Formula in Noncommutative Geometry and the Zeros of the Riemann Zeta Function,''
\emph{Selecta Mathematica}, vol. 5, no. 1, pp. 29--106, 1999.
DOI: \href{https://doi.org/10.1007/s000290050042}{10.1007/s000290050042}.

\bibitem{Sierra2007}
G. Sierra,
``$H=xp$ with Interaction and the Riemann Zeros,''
\emph{Nuclear Physics B}, vol. 776, pp. 327--364, 2007.
DOI: \href{https://doi.org/10.1016/j.nuclphysb.2007.03.049}{10.1016/j.nuclphysb.2007.03.049}.

\bibitem{deBruijn1950}
N. G. de Bruijn,
``The Roots of Trigonometric Integrals,''
\emph{Duke Mathematical Journal}, vol. 17, pp. 197--226, 1950.

\bibitem{Newman1976}
C. M. Newman,
``Fourier Transforms with Only Real Zeros,''
\emph{Proceedings of the American Mathematical Society}, vol. 61, no. 2, pp. 245--251, 1976.

\bibitem{RodgersTao2020}
B. Rodgers and T. Tao,
``The de Bruijn--Newman Constant Is Non-negative,''
\emph{Forum of Mathematics, Pi}, vol. 8, e6, 2020.
DOI: \href{https://doi.org/10.1017/fmp.2020.6}{10.1017/fmp.2020.6}.

\bibitem{deBranges1968}
L. de Branges,
\emph{Hilbert Spaces of Entire Functions},
Prentice-Hall, 1968.

\bibitem{KatzSarnak1999}
N. M. Katz and P. Sarnak,
\emph{Random Matrices, Frobenius Eigenvalues, and Monodromy},
American Mathematical Society, 1999.

\bibitem{Montgomery1973}
H. L. Montgomery,
``The Pair Correlation of Zeros of the Zeta Function,''
in \emph{Analytic Number Theory}, Proceedings of Symposia in Pure Mathematics, vol. 24, pp. 181--193, 1973.

\end{thebibliography}
